\documentclass[12pt,a4paper]{article}
\usepackage[utf8]{inputenc}
\usepackage{graphicx}
\usepackage{hyperref}
\usepackage{listings}
\usepackage{xcolor}
\usepackage{amsmath}
\usepackage{float}
\usepackage[margin=2.5cm]{geometry}

\definecolor{codegreen}{rgb}{0,0.6,0}
\definecolor{codegray}{rgb}{0.5,0.5,0.5}
\definecolor{codepurple}{rgb}{0.58,0,0.82}
\definecolor{backcolour}{rgb}{0.95,0.95,0.92}

\lstdefinestyle{mystyle}{
    backgroundcolor=\color{backcolour},   
    commentstyle=\color{codegreen},
    keywordstyle=\color{magenta},
    numberstyle=\tiny\color{codegray},
    stringstyle=\color{codepurple},
    basicstyle=\ttfamily\footnotesize,
    breakatwhitespace=false,         
    breaklines=true,                 
    captionpos=b,                    
    keepspaces=true,                 
    numbers=left,                    
    numbersep=5pt,                  
    showspaces=false,                
    showstringspaces=false,
    showtabs=false,                  
    tabsize=2,
    frame=single
}

\lstset{style=mystyle}

\title{SnakeGame Project Documentation}
\author{Developer: wartorle}
\date{\today}

\begin{document}

\maketitle
\begin{abstract}
Documentation for SnakeGame project - a classic "Snake" game implementation in C++ using Qt framework. Project developed as part of School-21 curriculum.
\end{abstract}

\tableofcontents

\section{Introduction}
\subsection{About the Project}
SnakeGame is a modern implementation of the classic arcade game "Snake". The project demonstrates the use of object-oriented programming in C++ combined with the Qt graphical framework.

\subsection{Key Features}
\begin{itemize}
    \item Full-featured game logic
    \item Qt-based graphical interface
    \item High score system
    \item Multi-level difficulty system
    \item Temporary speed boost
    \item Comprehensive documentation and testing
\end{itemize}

\section{Project Architecture}

\subsection{File Structure}
\begin{verbatim}
CPP3_BrickGame_v2.0-1/
├── src/
│   ├── brick_game/
│   │   └── snake/
│   │       ├── main.cpp
│   │       ├── mainwindow.h
│   │       └── mainwindow.cpp
│   ├── gui/
│   ├── tests/
│   ├── documentation.tex
│   └── hs_log.txt
└── Makefile
\end{verbatim}

\section{Classes and Structures}

\subsection{GameInfo\_t Structure}
Main structure containing all game data:

\begin{lstlisting}[language=C++,caption=GameInfo\_t structure]
struct GameInfo_t {
    int** field = nullptr;      // Game field
    int** snake = nullptr;      // Snake position
    int** next = nullptr;       // Next state
    int* tail_x = nullptr;      // Tail X coordinates
    int* tail_y = nullptr;      // Tail Y coordinates
    int x = 0;                  // Current X position
    int y = 0;                  // Current Y position
    int n_tail = 0;             // Tail length
    int pause = 0;              // Pause state
    int level = 1;              // Current level
    int high_score = 0;         // High score
    int score = 0;              // Current score
    int speed = 400;            // Game speed
    int key = 0;                // Current key
    int prev_key = 0;           // Previous key
    bool first_run = true;      // First run flag
};
\end{lstlisting}

\subsection{Enumerations}
\begin{lstlisting}[language=C++,caption=User actions]
enum UserAction_t { 
    Up,         // Move up
    Down,       // Move down
    Left,       // Move left
    Right,      // Move right
    Terminate,  // Terminate game
    Pause,      // Pause
    Action,     // Action (space)
    None        // No action
};

enum GameState { 
    MENU,       // Menu
    PLAY,       // Gameplay
    PAUSE,      // Pause
    GAME_OVER,  // Game over
    EXIT        // Exit
};
\end{lstlisting}

\subsection{MainWindow Class}
Main application class, inherited from QMainWindow.

\begin{lstlisting}[language=C++,caption=Main class fields]
class MainWindow : public QMainWindow {
    Q_OBJECT
    
    // Timers
    QTimer* boostTimer;     // Speed boost timer
    QTimer* gameTimer;      // Main game timer
    
    // Game states
    GameInfo_t currentState; // Current state for rendering
    GameInfo_t gameInfo;     // Main game state
    GameState state;         // Current game state
    
    // Apple position
    int apple_height = 0;
    int apple_width = 0;
};
\end{lstlisting}

\section{Main Functions}

\subsection{Memory Management}
\begin{lstlisting}[language=C++,caption=Memory management functions]
// Allocate memory for game field
static void allocate_memory_for_field(GameInfo_t* gameInfo);

// Allocate memory for next state
static void allocate_memory_for_next(GameInfo_t* gameInfo);

// Allocate memory for snake
static void allocate_memory_for_snake(GameInfo_t* gameInfo);

// Free memory
static void free_memory_field(GameInfo_t* gameInfo);
static void free_memory_next(GameInfo_t* gameInfo);
static void free_memory_snake(GameInfo_t* gameInfo);
static void free_tail_arrays(GameInfo_t* gameInfo);
\end{lstlisting}

\subsection{Game Logic}
\begin{lstlisting}[language=C++,caption=Game logic functions]
// Game initialization
void initialization();
void reset();

// Apple generation
void generate_apple_in_field(GameInfo_t* gameInfo, 
                           int* apple_height, int* apple_width);

// Input handling
bool fix_buttons(UserAction_t newDir, UserAction_t prevDir);
int get_key_code(UserAction_t action);
UserAction_t get_signal(int user_input);
void handle_action(UserAction_t action, GameInfo_t* gameInfo);

// Movement and collisions
int collision_check(const GameInfo_t& gameInfo);
void snake_movement(GameInfo_t* gameInfo);
void clear_apple(GameInfo_t* gameInfo);

// State updates
GameInfo_t updateCurrentState(const GameInfo_t& gameInfo);
void actual_level(GameInfo_t* gameInfo);

// High score system
int load_high_score(const QString& filename);
void save_high_score(const QString& filename, int score);
\end{lstlisting}

\subsection{Game States}
\begin{lstlisting}[language=C++,caption=Game state functions]
void game_state_menu();      // Menu state
void game_state_pause();     // Pause state
void game_state_play();      // Main gameplay
void game_state_game_over(); // Game over state
void updateGame();           // Game update
\end{lstlisting}

\subsection{Event Handlers}
\begin{lstlisting}[language=C++,caption=Event handlers]
// Key press handling
void keyPressEvent(QKeyEvent* event) override;

// Game rendering
void paintEvent(QPaintEvent* event) override;
\end{lstlisting}

\section{High Score System}

The project uses a simple text-based high score system. The high score is saved in the \texttt{hs\_log.txt} file in plain text format.

\subsection{High Score Functions}
\begin{lstlisting}[language=C++]
// Load high score from file
int load_high_score(const QString& filename);

// Save high score to file
void save_high_score(const QString& filename, int score);
\end{lstlisting}

\section{Level System}

Difficulty level automatically increases with score:

\begin{table}[H]
\centering
\begin{tabular}{|c|c|c|}
\hline
\textbf{Score} & \textbf{Level} & \textbf{Speed (ms)} \\
\hline
0-4 & 1 & 400 \\
5-9 & 2 & 350 \\
10-14 & 3 & 300 \\
15-19 & 4 & 250 \\
20-24 & 5 & 220 \\
25-29 & 6 & 200 \\
30-34 & 7 & 180 \\
35-39 & 8 & 160 \\
40-44 & 9 & 140 \\
45+ & 10 & 120 \\
\hline
\end{tabular}
\caption{Difficulty level system}
\end{table}

\section{Controls}

\subsection{Control Keys}
\begin{itemize}
    \item \textbf{Up Arrow} - Move up
    \item \textbf{Down Arrow} - Move down
    \item \textbf{Left Arrow} - Move left
    \item \textbf{Right Arrow} - Move right
    \item \textbf{Enter} - Pause/Resume
    \item \textbf{Space} - Temporary speed boost
    \item \textbf{Escape} - Exit game
\end{itemize}

\section{Building and Running}

\subsection{Project Compilation}
\begin{verbatim}
# Create build directory
mkdir build
cd build

# Generate Makefile
cmake ..

# Compile
make

# Run tests
make test

# Run application
./SnakeGame
\end{verbatim}

\subsection{Documentation Generation}
\begin{verbatim}
# Compile LaTeX documentation
latex documentation.tex
latex documentation.tex

# View DVI
xdvi documentation.dvi

# Convert to PDF
dvipdf documentation.dvi
\end{verbatim}

\section{Testing}

The project includes comprehensive testing using Google Test. Tests cover:

\begin{itemize}
    \item Snake movement in all directions
    \item Collision detection
    \item High score system
    \item Level progression
    \item User input handling
    \item All game states
\end{itemize}

\section{Conclusion}

SnakeGame is a full-featured implementation of the classic game with modern development approaches. The project demonstrates best practices in object-oriented programming, memory management, and cross-platform application development.

\end{document}